\chapter[LITERATURE REVIEW]{Literature Review}
\label{ch:literature-review}

\section{Introduction}
\label{sec:lr-introduction}

This chapter reviews the literature in seven parts, each addressing a question
raised by the one before it. Section~\ref{sec:lr-corrosion} establishes what
corrosion is, which forms it takes, and what an inspection method would have to
measure. Section~\ref{sec:lr-corrosion-detection} examines how corrosion is
detected in practice and what current methods fail to recover.
Section~\ref{sec:lr-modalities} considers what each sensing modality can
physically contribute. Section~\ref{sec:lr-fusion} reviews how modalities are
combined, how they are spatially registered, and whether the evidence shows an
added modality improving performance. Section~\ref{sec:lr-architectures} surveys
the learning architectures that consume these streams, and
Section~\ref{sec:lr-datasets} examines the data and evaluation practice on which
reported results rest. Section~\ref{sec:lr-gap} draws the strands together into
a statement of the research gap.

The review covers publications from 2021 onwards, retrieved from Scopus, Web of
Science, IEEE Xplore and ScienceDirect. Original journal research articles are
preferred throughout; conference papers and preprints are cited where they
report evidence unavailable elsewhere, and are identified as such. Review
articles are cited for context but do not contribute rows to the comparative
tables.

Much of the fusion literature reviewed in Section~\ref{sec:lr-fusion} is
directed at applications other than corrosion. This is deliberate. Multimodal
fusion matured elsewhere, and the registration techniques, fusion architectures
and evaluation practices developed there are what a corrosion application would
have to draw on. The distance between those
domains and this one is itself part of the argument developed in
Section~\ref{sec:lr-gap}.

\section{Corrosion}
\label{sec:lr-corrosion}

Corrosion proceeds by an electrochemical mechanism. An electrolyte in contact
with the metal permits ion transfer between anodic and cathodic regions; iron is
oxidized at the anode and the liberated electrons migrate to the cathode, where
they reduce dissolved oxygen and water. Where two dissimilar metals share an
electrolyte, the less noble becomes the anode and corrodes preferentially while
the more noble is protected \citep{Qin2026,Wojciechowski2025}. The substrates affected span
industrial practice as a whole: buried and above-ground pipelines
\citep{Akhlaghi2023,Alvani2026}, reinforcement embedded in concrete \citep{Yang2025}, stainless steel process equipment
\citep{Li2022,Wojciechowski2025}, ship hulls and marine components \citep{Wang2021,Arias2025}, and
structural steel under atmospheric exposure \citep{DeKerf2022}.

The forms corrosion takes are conventionally distinguished by morphology.
Uniform corrosion proceeds evenly across an exposed surface because anodic and
cathodic sites are microscopic and shift continuously, so no location remains
anodic long enough to develop a pit; the resulting loss is predictable and can
be expressed as a rate in millimeters per year, which permits a corrosion
allowance to be designed into the structure \citep{Yang2026a,Arias2025}. Localized forms
behave differently. In pitting, a break in the passive film establishes a small
anode within a large cathode, and the resulting current concentration drives
deep penetration at a single location; the pit interior acidifies as metal ions
hydrolyze and chloride migrates inward, so the attack sustains itself
\citep{Nguyen2026,Liu2026,Li2022}. It is the most common of the localized forms,
and \citet{Huang2026a} state that the variability of its condition makes pitting
considerably harder to monitor than uniform attack, which is the practical
counterpart of the mechanism just described: uniform loss can be sampled
anywhere, while a pit has to be found before it can be measured. Crevice corrosion arises by the same differential-aeration
mechanism in oxygen-starved gaps beneath gaskets, deposits and joints \citep{Nguyen2026}.
Galvanic corrosion follows from dissimilar-metal contact, its severity governed
by the anode-to-cathode area ratio, while intergranular attack proceeds along
chromium-depleted zones adjacent to grain-boundary carbides, leaving a surface
that appears sound while cohesion is lost \citep{Qin2026,Wojciechowski2025}. Stress corrosion cracking
requires a susceptible alloy, tensile stress and a specific corrosive species
acting together; cracks initiate at pits and propagate with little visible
product, so a ductile metal can fail in brittle fashion below its yield strength
\citep{Liu2026,Akhlaghi2023}. Microbially influenced corrosion is driven principally by
sulfate-reducing bacteria, whose biofilms establish a local chemistry
independent of the bulk water and are held responsible for a substantial share
of corrosion damage in marine service \citep{Liang2021}.

Corrosion is also classified by exposure environment, a separate axis along
which the morphology within each class varies with local conditions.
Atmospheric corrosion begins once relative humidity exceeds a critical value and
a thin electrolyte film forms, so the time of wetness governs the rate;
dissolved carbon dioxide and nitrogen oxides acidify the film and raise its
conductivity \citep{Arias2025,DeKerf2022}. In soil, uneven moisture, oxygen and chemistry
along a buried length render one stretch anodic relative to another, forming
macro-cells that span meters of pipe and favoring localized attack over uniform
loss \citep{Akhlaghi2023}. In oil and gas service, carbon dioxide dissolved in produced water
forms carbonic acid and attacks preferentially where the protective iron
carbonate scale breaks down, while hydrogen sulfide forms sulfide films and
supplies atomic hydrogen, raising the risk of sulfide stress cracking
\citep{Huang2026a,Alvani2026}.

The factors that govern all of these are broadly common: chloride
concentration, dissolved oxygen, temperature, relative humidity, pH,
atmospheric pollutants, surface condition, and the presence of microorganisms
\citep{Liu2026,Li2022,Liang2021}. Their consequences, however, are not equivalent.
Table~\ref{tab:corrosion} summarizes the corrosion studies reviewed, recording
for each the form addressed, the driving factors, the detection method where one
is applied, the severity measure reported, and the class of structure studied.

\begingroup
\scriptsize
\renewcommand{\arraystretch}{1.15}
\begin{xltabular}{\textwidth}{@{}l X X X l@{}}
\caption{Summary of corrosion studies reviewed.}\label{tab:corrosion}\\
\toprule
\textbf{Form} & \textbf{Driving factors} & \textbf{Detection method} & \textbf{Severity reported} & \textbf{Ref.} \\
\midrule
\endfirsthead
\multicolumn{5}{@{}l}{\footnotesize\itshape Table~\thetable\ (continued)}\\
\toprule
\textbf{Form} & \textbf{Driving factors} & \textbf{Detection method} & \textbf{Severity reported} & \textbf{Ref.} \\
\midrule
\endhead
\midrule
\multicolumn{5}{r@{}}{\footnotesize\itshape Continued on next page}\\
\endfoot
\bottomrule
\endlastfoot
Pitting & Hydrochloric acid & Multi-spectral imaging, thermocouple, PCA and FNN; 0.2829 mm depth error & Metal loss depth & \citep{Egodawela2025} \\
Pitting & Chloride, compositional inhomogeneity & --- & Pit depth to 1305.4 \si{\micro\meter} & \citep{Qin2026} \\
Pitting & Chloride, oxygen, iron-oxidizing bacteria & --- & --- & \citep{Nguyen2026} \\
Pitting & Chloride, temperature, stress & --- & Pit depth to 3.7 mm, width to 18 mm & \citep{Liu2026} \\
Pitting & Chloride, hypochlorite, pH & --- & --- & \citep{Li2022} \\
Pitting & Pressure, temperature, oxygen, chloride & --- & Pit depth to 10.86 \si{\micro\meter} & \citep{Wang2021} \\
Pitting & Soil properties, chloride, sulfate, pH & --- & Pit depth 0.41--13.44 mm & \citep{Akhlaghi2023} \\
Pitting & --- & Ultrasonic thickness, FSM and ANN & Over half the wall thickness & \citep{Huang2026a} \\
Pitting & Humidity, temperature, pollutants & Ultrasonic, RGB and LiDAR, cross-attention fusion; 98.44\% & Ordinal grade; modality ablation reported & \citep{Ding2025} \\
Pitting & Chloride, air, water, pollution & Active thermography with mHCD-2 & --- & \citep{Elizalde2023} \\
Pitting, intergranular & Anodic polarization, sodium sulfate & --- & --- & \citep{Wojciechowski2025} \\
Uniform & Chloride & Acoustic sensing with LBP, HOG and SVM; 98.14\% & --- & \citep{Yang2026a} \\
Uniform & Chloride, humidity & Pulsed eddy current with TMR sensor & Rebar 18 mm to 10.58 mm & \citep{Yang2025} \\
Atmospheric & Chloride, humidity, pollutants & VIS--NIR hyperspectral with MLP & Five ordinal classes & \citep{Arias2025} \\
Atmospheric & Chloride & SWIR hyperspectral with random forest; 3.27--13.37\% error & --- & \citep{DeKerf2022} \\
Microbial & Sulfate-reducing bacteria & --- & --- & \citep{Liang2021} \\
\end{xltabular}
\endgroup

Two observations follow from the table. Pitting is the form most frequently
addressed, appearing in 11 of the 25 studies reviewed across this section and
Section~\ref{sec:lr-corrosion-detection}, and the reason is a mismatch between
what the surface reveals and what the metal has lost: chloride attack drives
deep into the substrate while the surface presents only superficial rust
\citep{Nguyen2026,Akhlaghi2023}. Field measurements bear this out, with reported pit depths ranging
from 10.86 \si{\micro\meter} on marine stainless components \citep{Wang2021} to 13.44 mm
in buried transmission pipelines \citep{Akhlaghi2023}.

The second observation concerns measurement. Thirteen of the 25 studies never
quantify material loss at all, reporting electrochemical current or a corrosion
rate instead, quantities that describe how rapidly attack proceeds rather than
its physical extent \citep{Liang2021}. Of the 12 that do quantify damage, the measures
are mutually incompatible: pit depth \citep{Qin2026,Liu2026,Wang2021,Akhlaghi2023}, percentage of surface area
affected \citep{Alvani2026}, remaining wall thickness \citep{Yang2025}, or an ordinal grade
\citep{Arias2025,Bui2026,Ding2025}. \citet{Alvani2026} states the limitation explicitly, noting that
image-based severity assessment cannot measure material loss or corrosion depth.
No common basis for comparison therefore exists across the literature, and any
method claiming to recover true material loss must establish its own.

\section{Corrosion Inspection and Detection}
\label{sec:lr-corrosion-detection}

Detection of corrosion has traditionally relied on non-destructive testing.
Ultrasonic thickness measurement, pulsed eddy current, magnetic flux leakage,
radiographic testing and acoustic emission are all established in industrial
practice \citep{Huang2026a,Yang2026a,Yang2025}, and each returns a physical quantity rather than an
appearance. Their limitations are practical rather than physical: the equipment
is bulky and costly, deployment requires surface access and skilled operators,
and scalability to continuous monitoring of a large asset is correspondingly
limited \citep{Alvani2026}. At the scale of a large industrial asset this is
impractical \citep{Imran2025}.

Optical inspection offers contactless acquisition at a rate that non-destructive
testing cannot match, and the recent literature is dominated by learned
image-based methods. Three task formulations recur. Classification assigns a
label to an image or a patch, and has been applied to hull plates using AlexNet
with a sliding-window scan \citep{Yao2019} and to gas pipeline imagery using a VGG16
backbone with an XGBoost head \citep{Alvani2026}. Segmentation labels each pixel, and has
been applied to marine vessel hulls with a
U-Net, BiLSTM and DFE-NET combination reporting an F1-score of 0.8319 \citep{Yu2022},
and to industrial steel structures with a Bayesian HRNetV2 that additionally
returns per-pixel confidence \citep{Nash2022}. Detection returns bounding boxes and is
represented by YOLO-family models applied to reinforced concrete via GPR
B-scans \citep{Bui2026}. Classical image processing
retains a place: \citet{Imran2025} report 96.52\% segmentation accuracy using
filtering with a Chan--Vese active contour on ship structures, without a learned
model. What displaced such methods as the default was not accuracy but labor.
\citet{Yao2019} observe that grading corrosion through hand-crafted feature
extraction ties the method to the particular structure it was designed around
and carries a workload in manpower and cost that does not scale, whereas a
convolutional network derives the features from the data provided the samples
are there.

Reported accuracies across this literature are high, frequently exceeding 90\%
\citep{Alvani2026,Imran2025,Yao2019,Lavadiya2022}. The figures are not directly comparable, however, since the
task formulations differ, the datasets are private and locally collected, and
the severity measures are inconsistent. \citet{Nash2022} documents the
underlying difficulty: in the absence of a public corrosion dataset, models are
trained on data suited to immediate local conditions and fail to generalize
beyond them, producing substantial false positive and false negative rates when
deployed in the field. \citet{Alvani2026} names the sources of the variation
that defeats them: coatings, illumination, camera optics and surface background
all differ across operators and inspection campaigns, so the same corroded
surface does not present the same appearance twice.

The ceiling is recognized from inside this literature rather than only from
outside it. \citet{Imran2025} identify their own reliance on a single imaging
modality as the constraint on what their framework can adapt to, and name the
combination of optical, thermal and acoustic sensing as the route to corrosion
that one modality does not capture, assigning the thermal channel the specific
role of revealing the temperature variations that indicate loss beneath the
surface. Their third channel is sonar rather than a ranging sensor, which the
submerged setting of their application makes the natural choice. What the
passage establishes is not a particular sensor set but where the limit sits: in
what is sensed rather than in how it is processed.

A smaller group of studies combines more than one sensing modality.
\citet{Huang2026a} pair ultrasonic thickness measurement with field-signature
sensing and an artificial neural network on buried pipelines, and
\citet{Haitz2022} fuses industrial RGB imagery with laser triangulation range and
intensity on metal barrels, constructing a five-channel feature space.
Table~\ref{tab:corrosion-detection} summarizes the detection studies reviewed.

\begin{table}[htbp]
\centering
\caption{Summary of corrosion detection studies reviewed.}
\label{tab:corrosion-detection}
\scriptsize
\renewcommand{\arraystretch}{1.15}
\begin{tabularx}{\textwidth}{@{}X X X X l@{}}
\toprule
\textbf{Sensing modality} & \textbf{Method and performance} & \textbf{Severity reported} & \textbf{Structure} & \textbf{Ref.} \\
\midrule
RGB & VGG16 with squeeze-and-excitation and XGBoost head & 1.5\% to over 6.0\% of area & Gas transmission pipelines & \citep{Alvani2026} \\
RGB & Filtering with Chan--Vese active contour; 96.52\% segmentation accuracy & --- & Ships and naval vessels & \citep{Imran2025} \\
RGB & U-Net, BiLSTM and DFE-NET; F1 0.8319 & --- & Marine vessel hulls & \citep{Yu2022} \\
RGB & AlexNet with overlap-scanning sliding window; 98.9\% patch classification accuracy & --- & Hull structural plates & \citep{Yao2019} \\
GPR B-scan & YOLOv11 & Four ordinal levels & Reinforced concrete slabs & \citep{Bui2026} \\
RGB & Bayesian HRNetV2; F1 0.89 & --- & Steel industrial structures & \citep{Nash2022} \\
Cold thermography & Dynamic time warping & Metal loss up to 50\% of wall thickness & Mild carbon steel & \citep{Doshvarpassand2022} \\
Hyperspectral & Support vector machine; 93.2\% & --- & ASTM A572 structural steel & \citep{Lavadiya2022} \\
RGB and laser profilometry & Random forest on a 5D feature space; 0.97 & --- & Metal cylindrical barrels & \citep{Haitz2022} \\
\bottomrule
\end{tabularx}

\vspace{4pt}
\scriptsize Dashes denote information not reported by the study.
\end{table}

Three findings follow. First, multi-sensor fusion is rare in this literature.
Of the 25 studies reviewed across this section and
Section~\ref{sec:lr-corrosion}, 7 perform no detection study at all, and of the
18 that do, only 3 combine more than one sensing modality
\citep{Huang2026a,Ding2025,Haitz2022}, and 2 of those 3
combine more than one imaging or ranging sensor
\citep{Ding2025,Haitz2022}. A fourth,
\citet{Egodawela2025}, pairs multispectral imaging with a contact thermocouple,
which is a point probe rather than an imaging channel and is counted here as
single-sensor imaging.

Second, one study already does what the rest do not, and it warrants close
attention. \citet{Ding2025} is a conference paper and sits outside the peer
reviewed journal corpus this review tabulates; it is therefore cited as evidence
that a thing has been attempted, never as a comparator and never as support for
the gap this chapter argues. It combines ultrasonic guided waves,
visible imagery and LiDAR point clouds on aerial ground wires, and fuses them
inside the model: three parallel encoders, intra-modal self-attention within each
stream and multi-head cross-attention between them, trained jointly for corrosion
classification and severity grading on a purpose-built registered dataset. It
also reports a leave-one-out modality ablation, the only one in the corrosion
literature reviewed here, finding that removing the ultrasonic channel costs
9.2\% accuracy, the visible channel 3.6\% and the point cloud 2.9\%. Learned
trimodal fusion for a corrosion decision is therefore not unattempted. What
constrains it is the sensing: the ultrasonic transducers are mounted directly on
the conductor surface, so the system as a whole is contact-based and inherits
the access limitation that optical inspection exists to remove. Severity is
reported as an ordinal class rather than a physical measurement. The two
remaining multi-sensor systems are applied to buried pipelines
\citep{Huang2026a} and to a barrel on a rotary-table inspection rig
\citep{Haitz2022}: industrial assets in both cases, and in the second a bench
rather than a structure in service.

Third, none of them recovers a depth measurement by fusion. Depth is recovered
in this literature, but by single modalities and under controlled conditions.
\citet{Doshvarpassand2022} matches thermal contrast against an analytic model for
losses up to half the wall thickness under cryogenic stimulation, and \citet{Egodawela2025} regresses metal
loss from multi-spectral imaging with a reported depth error of 0.2829 mm,
demonstrating that depth need not come from a geometric sensor. Neither, however, operates at the standoff or under
the motion that field inspection of a large structure imposes.

A gap therefore exists for contactless multi-sensor fusion capable of recovering
pitting depth in the field. No journal study reviewed here fuses more than two
sensing modalities for a corrosion decision, and none recovers a depth
measurement from a fused representation at all. The single system that carries a
ranging channel alongside imagery for corrosion requires surface-mounted
transducers on a third \citep{Ding2025}, and it reports a grade rather than a
measurement.

\section{Sensing Modalities for Corrosion Assessment}
\label{sec:lr-modalities}

Section~\ref{sec:lr-corrosion-detection} established that current detection
methods recover appearance and affected area but rarely depth. This section examines what each
candidate sensing modality can physically measure, since the case for combining
them rests on their measuring different things rather than on the number of
sensors carried.

\subsection{Visible Imaging}

A color camera records light reflected from a surface, encoding the color and
texture that distinguish oxidized metal from sound metal. Its advantages are
practical: sensors are inexpensive and available at high resolution, acquisition
is fast and contactless, and a mature ecosystem of pre-trained architectures
transfers readily to the task \citep{Yao2019,Nash2022}. The great majority of the
corrosion detection literature reviewed in
Section~\ref{sec:lr-corrosion-detection} rests on this modality alone
\citep{Alvani2026,Imran2025,Yu2022}.

Its limitations are physical rather than algorithmic, and three recur.
Corrosion products, coating discoloration, surface staining and shadow present
similar color and texture, so a model trained on appearance cannot reliably
separate them; \citet{Lavadiya2022} identifies this visual ambiguity explicitly, and
\citet{Liu2026} attribute the residual segmentation error of loss functions
operating on color and texture alone to precisely these confounders. Second, a
single image carries no depth information \citep{Cheng2025}, so severity can be
expressed only as the proportion of surface area affected; \citet{Alvani2026} states
the consequence directly, noting that image-based severity cannot measure
material loss or corrosion depth. \citet{Egodawela2025} puts the difficulty in
its general form, that grading metal loss from imaging requires depth to be
inferred from two-dimensional data rather than measured, which is the same
distinction Section~\ref{sec:lr-geometric} draws between a depth that is
reconstructed and a depth that is ranged. Third, corrosion progressing beneath an intact
coating produces no reflected-light signature at all \citep{Lim2021}.

\subsection{Infrared Thermography}

Thermography measures emitted infrared radiation, and therefore surface
temperature. \citet{Egodawela2025} judges it the most successful imaging
technique presently available for quantifying metal loss, on the grounds that
visible imaging is confined to surface-level defects while terahertz imaging
struggles against the reflectivity of metals, leaving heat conduction as the
property that carries subsurface information to a camera. Its relevance to corrosion is indirect but powerful: subsurface
material loss alters the local heat capacity and conduction path, so a defect
below an intact surface produces a thermal contrast even where it produces no
visual one. The governing physics is diffusion. Heat propagates into the
material according to the Fourier equation, at a rate set by the thermal
diffusivity $\alpha = k/\rho C$, and a subsurface cavity disrupts that
propagation, so the disruption surfaces as a temperature anomaly over the
defect \citep{Doshvarpassand2022}. Passive thermography relies on ambient
thermal gradients, while active thermography applies controlled excitation and
observes the transient response \citep{Elizalde2023}. The same diffusion that
carries the signal outward is what limits how deep it can be read from, which
the limitations below return to.

The evidence that this recovers what visible imaging cannot is direct.
\citet{Lim2021} fuses visible and active infrared images with a Faster R-CNN on
steel bridges and detect subsurface corrosion invisible in the color channel,
using halogen lamps as the heat source. \citet{Doshvarpassand2022} estimate defect depth from cold-stimulated
thermography on flat-bottom holes drilled into the hidden face of a mild steel
plate, and report the potential to predict metal loss of up to 50\% of wall
thickness. \citet{Elizalde2023} records corroded galvanized and stainless steel sheets
peaking at 63.0 and 65.3\,\si{\degreeCelsius} against 69.7 and
70.6\,\si{\degreeCelsius} for sound sheets, a separation of roughly 5 to
7\,\si{\degreeCelsius} in which the corroded surface is the cooler one, and
notes that this sign is the reverse of what the thermography literature
generally reports for hidden defects.

Two limitations qualify this, and the first is more
severe than it is usually presented. Every result above was obtained under
deliberate thermal excitation: \citet{Lim2021} uses halogen lamps,
\citet{Elizalde2023} two quartz-tube heaters totalling 2000\,W positioned
0.70\,m from the specimen with the camera at 1\,m, and
\citet{Doshvarpassand2022} an aerosol cryogen delivered from a rail-guided
carrier passing over the surface, an inversion of the usual arrangement that its
authors note is rarely reported against the prevailing use of heating. No reviewed study recovers subsurface
corrosion in steel passively. \citet{Yamashita2025} does work passively, but on
asphalt pavement over a concrete deck, at night to avoid solar loading, and reports that damage
progressing in the depth direction produces little surface temperature change.
The excitation arrangement therefore does not scale to inspection standoff on
the evidence available, and whether a passive regime can substitute is an open
question rather than an established capability. Excitation also changes what the
measurement is. In the active regime the camera records the temporal as well as
the spatial evolution of surface temperature, from the moment of stimulation
until the surface settles back to ambient \citep{Doshvarpassand2022}, so the
thermal channel is a sequence keyed to a stimulation event rather than a frame
that can be taken alongside a photograph. Bringing it into correspondence with
the other modalities is therefore a question of timing as well as of geometry,
and Chapter~\ref{ch:methodology} has to answer both. The second limitation follows from the
diffusion mechanism itself: thermal contrast attenuates with defect depth, since
the signal from a deeper cavity reaches the surface later and weaker.
\citet{Doshvarpassand2022}
reports accuracy falling as the defect aspect ratio falls, and notes from the
wider literature that depth prediction beyond half the wall thickness has barely
been reported. Beyond these two, reflective or transparent surfaces can produce
inaccurate thermal readings \citep{Tran2024}.

\subsection{Geometric Sensing}
\label{sec:lr-geometric}

Geometric sensors measure surface shape directly, and therefore measure the
quantity that corrosion severity is defined by. Several instruments serve this
role and they are not interchangeable. Laser triangulation and laser
profilometry recover fine surface relief at short, controlled standoff:
\citet{Haitz2022} adds triangulation range and intensity as the fourth and fifth
channels of a feature space built on a rotary table, recovering fine relief at a
standoff the rig holds fixed. LiDAR measures range by time of flight and operates at meters rather
than millimeters, tolerating platform motion and standoff variation that
profilometry does not. \citet{Maru2025} fuses LiDAR points with imagery at the
pixel level to assess section loss in steel structural components, and
\citet{PimenteldeFigueiredo2025} builds semantic--geometric maps of industrial environments by
combining point-cloud localization with vision-based segmentation.
\citet{Gunatilake2021} combine stereo vision with laser profiling to map internal
pipeline defects.

A distinction matters here that the literature does not always make. Depth
recovered by photogrammetric reconstruction or structure-from-motion is
inferred from imagery rather than measured, and inherits the failure modes of
the color channel it derives from; depth measured by a ranging instrument does
not \citep{Tang2025}. A separate caution concerns access rather than
inference. \citet{Ding2025} combines LiDAR and color imagery, both genuinely
measured and genuinely fused, with ultrasonic transducers mounted directly on
the conductor surface. The geometric and visible channels are contactless; the
system is not, and it is the contact channel that carries most of the
performance.

The limitations of geometric sensing are resolution and registration. Spatial
resolution must be sufficient to resolve a pit mouth measured in millimeters or
smaller \citep{Qin2026,Wang2021}, and sparse point clouds also limit object
detection \citep{Tran2024}. One reviewed study bears directly on whether a
ranging sensor meets that requirement at inspection standoff.
\citet{Maru2025} places a terrestrial laser scanner and a camera roughly four
metres from 500\,mm steel plates carrying acid-etched losses of 0.5, 1, 2 and
3\,mm, and recovers depth with a mean error of 7.8\%. Performance is depth
dependent in the expected direction: losses of 2\,mm and deeper are recovered
closely, while the 0.5\,mm case carries 17.2\% error, and estimates are biased
low at every depth. This establishes that depth recovery by a ranging sensor at
metre standoff is feasible for millimetre-scale losses and marginal below
roughly one millimetre, on prepared specimens without corrosion product. The
second limitation is that a point cloud must be brought into correspondence with
the image data before joint inference is possible, a problem
Section~\ref{sec:lr-fusion} treats in detail.

\subsection{Spectral Imaging}

Hyperspectral and multispectral imaging sample reflectance across many narrow
bands, separating materials by spectral signature rather than by apparent color.
Applied to corrosion, this addresses the ambiguity problem directly.
\citet{Lavadiya2022} demonstrates that hyperspectral imaging eliminates visual
ambiguity between corroded and coated surfaces and identifies corrosion sources.
\citet{Arias2025} grade marine steel into five severity classes from VIS--NIR
reflectance with a multilayer perceptron, and \citet{DeKerf2022} quantify corrosion
minerals in carbon steel from SWIR bands with a random forest, reporting errors
between 3.27\% and 13.37\%. \citet{Egodawela2025} estimates metal loss depth from
multi-spectral imagery with a reported error of 0.2829 mm.

Spectral imaging is therefore a credible alternative to thermography as the
second modality, and the reason it is not adopted here is one of deployment
rather than capability. A hyperspectral line scanner requires controlled
illumination and stable standoff to produce calibrated reflectance, conditions
that a fixed laboratory arrangement satisfies and a moving inspection platform
does not. Thermography, by contrast, responds to subsurface material loss rather
than surface chemistry, which is the quantity pitting demands, and imposes no
illumination requirement. Table~\ref{tab:modalities} summarizes the comparison.

\begin{table}[htbp]
\centering
\caption{Sensing modalities considered for corrosion assessment.}
\label{tab:modalities}
\scriptsize
\renewcommand{\arraystretch}{1.2}
\begin{tabularx}{\textwidth}{@{}l X X X l@{}}
\toprule
\textbf{Modality} & \textbf{Quantity measured} & \textbf{Contribution} & \textbf{Limitation} & \textbf{Ref.} \\
\midrule
Visible & Reflected light: color and texture & Mature, inexpensive, high resolution; identifies corroded appearance & Confuses rust with stain, coating and shadow; no depth; blind beneath coatings & \citep{Alvani2026,Nash2022,Lavadiya2022,Lim2021} \\
Thermography & Emitted infrared radiation, i.e.\ surface temperature & Reveals subsurface loss beneath an intact surface; supports depth estimation & Requires excitation for active use; contrast attenuates with depth & \citep{Lim2021,Doshvarpassand2022,Elizalde2023} \\
Laser profilometry & Surface relief at short standoff & Fine depth and volume of individual pits & Fixed standoff, bench-mounted; poor tolerance of platform motion & \citep{Haitz2022} \\
LiDAR & Range by time of flight & Geometry at meter standoff; tolerates motion; measures rather than infers depth & Resolution limits detection of small pits; requires registration to imagery & \citep{Maru2025,PimenteldeFigueiredo2025} \\
Hyperspectral & Reflectance across narrow bands & Separates corroded from coated surfaces by spectral signature & Needs controlled illumination and stable standoff; high data volume & \citep{Lavadiya2022,Arias2025,DeKerf2022} \\
\bottomrule
\end{tabularx}
\end{table}

No single modality recovers everything corrosion assessment requires. Visible
imaging localizes the affected region but cannot say how much metal is gone;
thermography reaches beneath the surface but needs excitation and loses
sensitivity with depth; geometric sensing measures loss directly but at a
resolution and registration cost; spectral imaging resolves material ambiguity
but under conditions field inspection does not provide. The complementarity is
the argument for combining them, and Section~\ref{sec:lr-fusion} examines how
that combination has been attempted.

\section{Multimodal Sensor Fusion}
\label{sec:lr-fusion}

Section~\ref{sec:lr-modalities} established that the candidate modalities
measure different physical quantities and fail in different ways \citep{Tran2024}. This section
examines how the literature combines them: at what stage the streams are joined,
how they are brought into spatial correspondence, and what evidence exists that
the additional modality improves performance.

One point of terminology has to be settled first, because the word is used in
two senses that do not overlap. \citet{Yao2026} surveys multimodal
computer vision for corrosion assessment, and the modalities it ranges over are
characterization techniques at different scales, from atomic force and scanning
electron microscopy through optical and multispectral imaging, combined to
understand corrosion behavior. Throughout this chapter multimodal means
something narrower: two or more sensors of different physical type, mounted
together and directed at the same surface during an inspection, whose outputs
are brought into a common frame and combined for a decision about that surface.
A survey of the first sense does not speak to the absence argued in the second.

\subsection{Fusion Architectures}

Fusion is conventionally classified by the stage at which streams are combined.
Early fusion concatenates raw or lightly processed channels at the input, so a
single backbone learns across modalities from the outset. \citet{Cao2026} builds
a five-channel input from color, inferred depth and thermal data;
\citet{Nooralishahi2022} aligns color, thermal and active-stereo depth at a fixed step to produce a
point cloud carrying both color and temperature; and \citet{Maru2025} projects
LiDAR points onto image pixels before feature extraction. Early fusion is
economical, since one backbone is trained rather than several, but it requires
the streams to be accurately registered beforehand, and \citet{Cao2026}
identifies exactly this as its limitation.

Intermediate fusion, in which separate encoders extract modality-specific
features that are combined before the detection head, dominates the reviewed
literature. Implementations differ mainly in the combination mechanism.
Cross-attention is the most common: \citet{Wang2026b} weights each stream by how
occluded it is, \citet{Jiang2024} pairs unimodal and cross-modal attention across
CNN and Transformer branches, and \citet{Cheng2025} uses LiDAR depth to guide image
features into a shared bird's-eye representation. Others fuse at multiple
pyramid levels \citep{Tang2026} or progressively across a branch tree that
discards redundant information \citep{Liu2025a}.

Late fusion combines decisions rather than features, joining the streams after
each has been processed independently \citep{Lin2022,Shahsavarani2024,Tran2024}. It tolerates sensor
dropout better but forgoes cross-modal interaction during learning.

The stage is not an arbitrary choice, and one study tests it directly.
\citet{Yang2023} compares early, intermediate and late fusion of color and
thermal imagery on the same facade defect data, reporting that the intermediate
variant achieves high accuracy with fewer weights while the early variant
preserves raw detail at the cost of doubled parameters. This is the only
fusion-stage comparison in the reviewed set.

\subsection{Registration and Calibration}

Fusion presupposes that measurements from different sensors refer to the same
physical point, and establishing that correspondence is the most frequently
reported obstacle in this literature. Nine of the tabulated studies name
alignment, calibration or registration failure among their stated limitations
\citep{Cheng2025,Liu2025,Zhu2026,Xu2026,Cao2026,Nooralishahi2022,Shahsavarani2024,Bin2023,Yang2024},
and the failure modes are specific: \citet{Zhu2026} reports breakdown under
sensor offset and motion blur, \citet{Xu2026} under large offsets, and
\citet{Nooralishahi2022} holds alignment only at the single standoff for which it was
calibrated.

The difficulty is compounded for thermal sensors, because the printed
chessboard used for conventional camera calibration has no thermal contrast and
is therefore invisible to an infrared detector. Two approaches follow. The first
is targetless calibration, estimating extrinsics from scene structure rather
than a known object: \citet{Yang2024} matches three-dimensional edges to image
lines to calibrate thermal, RGB and LiDAR sensors without any target, and
\citet{Tran2024} achieves the same for a 360-degree rig while the platform is
moving. The second is to design the network to tolerate misalignment:
\citet{Wang2024} reports salient object detection on unaligned and loosely aligned
pairs, and \citet{Liu2025} learns correction offsets internally, though it notes
that correcting alignment did not by itself raise detection accuracy.
\citet{Li2024c} performs coarse-to-fine registration and merges detections at the
head. Five studies resolve registration internally in one of these ways
\citep{Yang2024,Tran2024,Wang2024,Liu2025,Li2024c}. Two of them,
\citet{Yang2024} and \citet{Tran2024}, do so across exactly the three
modalities of interest here, which establishes that the correspondence is
recoverable; both, however, treat calibration as the end product and neither
carries the aligned streams into a detection model. Registration for this
sensor set is therefore a solved problem in robotics and an unused one in
inspection.

\subsection{Evidence of Modality Contribution}

Fusion is reported to outperform single-sensor baselines throughout this
literature, but the comparison is usually between a fused model and one
unimodal model rather than across the sensor set, so the contribution of an
individual modality is rarely isolated, though the exceptions are more numerous
than the field's own framing suggests. \citet{Hou2024}, an infrared--visible
detection study drawn on here for its ablation design rather than as a corrosion
comparator, trains the same detector on each modality separately and on their
fusion. \citet{Haitz2022} compares
seven channel combinations and finds the differences between them not highly
significant. \citet{Maru2025} compares the fused input against color alone and
reports a 15.1\% gain in intersection over union, though without a
geometry-only arm and without separating range from laser intensity.
\citet{Rashidi2026} runs six configurations and concludes that no single
modality suffices. \citet{Ding2025}, discussed in
Section~\ref{sec:lr-corrosion-detection}, reports the only leave-one-out
ablation on a corrosion task. What none of them isolates is the contribution of
a ranging channel by itself.

Within corrosion specifically the evidence is thinner, and it is not uniformly
favorable. \citet{Haitz2022} evaluates a five-channel feature space combining
color with laser triangulation range and intensity against a three-channel color
baseline, and reports the outcome as only slightly superior, noting that the
margins are not highly significant and closing with the question of whether the
active sensor is worth the fusion complexity it requires. That question is the
one Objective~3 of this research is designed to answer, and it has not been
answered since.

The counter-evidence must be stated alongside it. \citet{Wu2025} fuses color and
LiDAR with an inferred near-infrared channel for pedestrian detection and reports
the fused model weaker than single-sensor alternatives in both simple and
difficult scenes, and \citet{Liu2025} finds that correcting cross-modal
alignment did not by itself raise detection accuracy. Adding a modality is
therefore not automatically beneficial, which is precisely why the contribution
has to be measured rather than assumed.

Against this, \citet{Lim2021} fuses color with active infrared and detects
subsurface corrosion that the color channel cannot represent at all. This
establishes that a second modality can contribute something categorical rather
than incremental, although the study does not quantify the contribution as a
performance delta.

Table~\ref{tab:fusion} summarizes the fusion studies reviewed.

\begingroup
\scriptsize
\renewcommand{\arraystretch}{1.1}
\begin{xltabular}{\textwidth}{@{}p{2.5cm} X X l@{}}
\caption{Summary of multimodal fusion studies reviewed. Modalities in bold denote the
combination proposed in this research; inferred channels are those derived from imagery
rather than measured by a dedicated sensor.}\label{tab:fusion}\\
\toprule
\textbf{Modalities} & \textbf{Method and fusion stage} & \textbf{Application and reported limitation} & \textbf{Ref.} \\
\midrule
\endfirsthead
\multicolumn{4}{@{}l}{\footnotesize\itshape Table~\thetable\ (continued)}\\
\toprule
\textbf{Modalities} & \textbf{Method and fusion stage} & \textbf{Application and reported limitation} & \textbf{Ref.} \\
\midrule
\endhead
\midrule
\multicolumn{4}{r@{}}{\footnotesize\itshape Continued on next page}\\
\endfoot
\bottomrule
\endlastfoot
RGB and thermal & SeAFusion with YOLOv7, intermediate & Vehicles and pedestrians. Thresholds set by hand; single dataset & \citep{Byeon2026} \\
RGB and depth scan & WideResNet-50 and DGCNN branches, gated intermediate & Defects in 3D shapes. Cannot run on one sensor; synthetic and lab data & \citep{Asad2026} \\
RGB and laser profiling & Color projected onto scanned points & Pipeline internal defects. Errors on dark surfaces; roll about the pipe axis undetectable & \citep{Gunatilake2021} \\
RGB and LiDAR & Swin-Tiny with LiDAR guidance, bird's-eye view & 3D object detection. Heavy; requires exact calibration; urban only & \citep{Cheng2025} \\
RGB and LiDAR & ResNet or ConvNeXt with FSD branch, intermediate & 3D object detection. Handles only two modalities & \citep{Zhang2026d} \\
RGB and LiDAR & ResNet-50 with U-Net, point-to-pixel early fusion; fused vs RGB-only compared & Section loss in steel components at 4 m standoff. Laboratory only; no real corrosion texture & \citep{Maru2025} \\
RGB and LiDAR & U-Net with ResNet-34, late projection onto the point cloud & Corrosion on metal structures. Noisy; pose drift; needs good initialization & \citep{PimenteldeFigueiredo2025} \\
RGB and LiDAR & R3LIVE, color projected onto points & Crop growth estimation. Weak on chlorophyll; no yield estimate & \citep{Tang2025} \\
RGB and thermal & CNN and Transformer branches, intermediate and late & Vehicles and pedestrians. Doubles inference time and parameter count & \citep{Jiang2024} \\
RGB and thermal & Dynamic interaction module, intermediate & Salient objects. Structure and efficiency unresolved & \citep{Gao2026} \\
RGB and thermal & DAFM two-way cross-attention, intermediate & Pedestrians and vehicles. Too heavy for low-power hardware; fails on weak thermal & \citep{Wang2026b} \\
RGB and thermal & AFFNet with LFCMF at pyramid levels & Vehicles in drone imagery. Misses partly visible objects & \citep{Tang2026} \\
RGB and thermal & Tree-structured multibranch progressive, intermediate & Aerial targets. Weak on small objects & \citep{Liu2025a} \\
RGB and thermal & CMAM and THAM alignment, intermediate and head & Vehicles and people. Corrects alignment without raising accuracy & \citep{Liu2025} \\
RGB and thermal & YOLOv3 and SSD, late & Pedestrians and vehicles. Training pairs matched by hand & \citep{Lin2022} \\
RGB and thermal & Separate encoder branches, intermediate & Detection and segmentation. Loses texture where sensors overlap & \citep{Lin2026} \\
RGB and thermal & ResNet-50 backbones, intermediate & Pedestrians, vehicles, traffic lights. Breaks under sensor offset or motion blur & \citep{Zhu2026} \\
RGB and thermal & ResNet branches, intermediate & Objects in complex scenes. Fails when RGB is too dark for region masks & \citep{Xiao2026} \\
RGB and thermal & Spectral-Enhanced Transformer, intermediate & Surveillance and driving imagery. Breaks under large offset; too slow for live use & \citep{Xu2026} \\
RGB and thermal & Dense residual backbone, spatial and frequency branches & Detection and segmentation. Large and slow & \citep{Chen2026a} \\
RGB and thermal & Transformer and reconstructed convolution, intermediate & Objects in low light. Too complex for edge devices & \citep{An2026} \\
RGB and thermal & Diffusion backbone with twin guidance & Scene understanding. Slow; untested on video & \citep{Zhou2026} \\
RGB and active thermal & Faster R-CNN, input-level fusion & \textbf{Corrosion on steel bridges}. Requires halogen lamps as heat source & \citep{Lim2021} \\
RGB and thermal & Shift and scale estimation, early & Image registration. Breaks under motion blur; cannot fix viewpoint change & \citep{Bin2023} \\
RGB and thermal & VGG-like encoder and graph network, late & Defect mapping on structures. Accuracy depends on the sensor pair & \citep{Shahsavarani2024} \\
RGB and thermal & SACNet two-way attention, intermediate & Salient objects. Fails on hollow objects and under strong light & \citep{Wang2024} \\
RGB and thermal & Modified U-Net, early, intermediate and late compared & Loose and missing tiles. Long training; doubled weights in some variants & \citep{Yang2023} \\
RGB and thermal & Coarse-to-fine alignment, boxes merged at head & Industrial part defects. Struggles on irregular shapes & \citep{Li2024c} \\
Thermal and GPR & MobileNetV3 branches, bidirectional cross-modal attention and gated fusion, intermediate & Delamination in concrete decks. Cannot produce an output where one modality lacks coverage & \citep{Rashidi2026} \\
Thermal and laser profiling & Rule-based link of temperature to surface shape; authors' own term & Damage on concrete bridge decks. Passive, so night imaging; unusable for days after rain & \citep{Yamashita2025} \\
RGB, LiDAR and inferred NIR & YOLOv5 and PointPillars, early and late & Pedestrians in low light. Weaker than alternatives in simple and hard scenes & \citep{Wu2025} \\
RGB, inferred depth and thermal & GLS-YOLOv8n, five-channel early fusion & Mangoes in orchards. Requires the three views aligned by hand & \citep{Cao2026} \\
RGB, thermal and active-stereo depth & Fixed-step alignment, early & Inspection of industrial components. Alignment holds at one standoff only & \citep{Nooralishahi2022} \\
RGB, thermal and NIR & ConvNeXt backbone, intermediate & Image fusion for detection. Raises contrast but loses fine gray detail & \citep{Laidouni2025} \\
RGB, thermal and multispectral & Stitched aerial maps compared layer by layer & Defects on building facades. Flights restricted to dry season by the platform; thermal and multispectral resolution too coarse for microcracks & \citep{Guedes2025} \\
\textbf{RGB, thermal and LiDAR} & Swin Transformers, bird's-eye view & 3D road targets. Overfits at close range; streaks where views overlap & \citep{Tran2025} \\
\textbf{RGB, thermal and LiDAR} & 3D edges matched to image lines & Sensor calibration for robots. Needs clear edges; blurred thermal outlines degrade it & \citep{Yang2024} \\
\textbf{RGB, thermal and LiDAR} & SuperPoint and MobileNetV3 branches, late & 360-degree sensor calibration. Fails on heat-absorbing or scattering surfaces & \citep{Tran2024} \\
\end{xltabular}
\endgroup

Three findings follow from the table. First, multimodal fusion is mature in
autonomous driving, remote sensing and surveillance, and largely absent from
defect detection. Color paired with thermal is the default configuration,
adopted in 21 of the 38 studies tabulated once active thermography is included,
which leaves depth information absent across most of the field.

Second, only 3 of the 38 sense color, thermal and LiDAR concurrently
\citep{Tran2025,Yang2024,Tran2024}, and none of the three performs damage assessment: two are
calibration studies \citep{Yang2024,Tran2024} and the third detects road targets for
autonomous driving \citep{Tran2025}. Where a third channel does appear in an
inspection context it is either inferred from imagery \citep{Cao2026} or
measured by a short-range active-stereo sensor whose alignment holds at one
standoff only \citep{Nooralishahi2022}, rather than by a ranging instrument
that tolerates the standoff and motion of field inspection.

Third, only two of the tabulated studies address corrosion at all.
\citet{Lim2021} fuses color with active thermal imagery on steel bridges, and
\citet{PimenteldeFigueiredo2025} combines color with LiDAR to map corrosion in industrial
environments. Neither carries all three modalities, and no study in this set
both senses three modalities and learns a joint representation from them for a
corrosion decision.

\section{Learning Architectures for Detection and Fusion}
\label{sec:lr-architectures}

Section~\ref{sec:lr-fusion} examined how streams are combined. That combination
is always hosted inside a detection architecture, and the architecture bounds
what a fused representation can express. This section reviews the detectors,
backbones and analytical tools the reviewed literature employs, in order to
establish which of them are available to this research and what they have so far
been asked to do.

\subsection{Detection Architectures}

Detection architectures divide into two families. Two-stage detectors generate
region proposals and then classify them, and the Faster R-CNN family remains in
active use for this purpose \citep{Wu2026,Zheng2023,Gong2022}. \citet{Wu2026}
reduces false detections through interval normalization weighting while still
missing occluded objects.
One-stage detectors regress boxes directly and dominate recent work
\citep{Han2022,Tao2021,Sun2023,Liu2026a,Wang2026a,Sun2022,Li2024b,Morita2024,Zhao2026}. \citet{Srivastava2021} compares the two families on identical
data and finds the one-stage YOLOv3 both fastest and strongest overall, but
weakest on small targets and under scale change. Anchor-free variants built on
FCOS occupy a third position \citep{Chen2025,Zhang2024}, and transformer and foundation-model
adaptations are beginning to appear \citep{Zhang2026c,Mbietieu2025}.

Two properties of this literature bear directly on the present work. The first
is that small objects remain the standing limitation: 6 of the tabulated
studies are directed at small or multi-scale targets
\citep{Tao2021,Chen2025,Jiang2025,Zhang2024,Wang2026a,Han2022}, and the problem
is reported as unresolved even where it is the stated objective:
\citet{Tao2021} misses objects against cluttered backgrounds despite being built
for small flying objects. The same pattern appears in the family comparison reported above: the
detector that wins overall is the one that loses most on small targets
\citep{Srivastava2021}. Pitting corrosion is a small-object problem, and a
detector inherits this weakness whatever it is fed.
The second is that detection is consistently stronger than measurement. \citet{Li2024b}
sorts bridge-bolt corrosion into three severity classes but assigns those classes
by hand, \citet{Zhao2026} notes that a bounding box cannot represent the outline of the
region it encloses, and \citet{Sun2022} localizes boiler wall defects without
quantifying them. Where measurement is attempted from imagery alone it degrades
sharply, as \citet{Su2024} shows in reporting 99.75\% accuracy for crack detection
alongside a coefficient of determination of 0.764 for the corresponding width
regression.

The single-sensor studies are collected in Table~\ref{tab:detection}. Of the 28
detection studies examined, 21 employ a single sensor, and 4 of those sense
something other than visible light \citep{Morita2024,Zhao2026,He2026,Gong2022}.
The remaining 7 are multimodal and are tabulated with the fusion literature in
Table~\ref{tab:fusion}.

\begingroup
\scriptsize
\renewcommand{\arraystretch}{1.1}
\begin{xltabular}{\textwidth}{@{}X l X X l@{}}
\caption{Summary of single-sensor detection studies reviewed. Applications in bold denote
metallic defect or corrosion targets.}\label{tab:detection}\\
\toprule
\textbf{Method} & \textbf{Data} & \textbf{Application} & \textbf{Reported limitation} & \textbf{Ref.} \\
\midrule
\endfirsthead
\multicolumn{5}{@{}l}{\footnotesize\itshape Table~\thetable\ (continued)}\\
\toprule
\textbf{Method} & \textbf{Data} & \textbf{Application} & \textbf{Reported limitation} & \textbf{Ref.} \\
\midrule
\endhead
\midrule
\multicolumn{5}{r@{}}{\footnotesize\itshape Continued on next page}\\
\endfoot
\bottomrule
\endlastfoot
YOLOv3, SSD and Faster R-CNN compared & RGB & Everyday object detection & Struggles with small targets and scale change & \citep{Srivastava2021} \\
MVNet on SSD with multi-receptive field modules & RGB & Aircraft, cars and ships in remote imagery & 38.74 billion operations per image & \citep{Han2022} \\
Faster R-CNN with interval normalization weighting & RGB & Objects in everyday scenes & Misses occluded objects & \citep{Wu2026} \\
LSL-Net, a lightweight YOLOv4-tiny variant & RGB & Small low-altitude flying objects & Misses objects against cluttered backgrounds & \citep{Tao2021} \\
FRLI-Net, FCOS with feature reconstruction & RGB & Tiny objects and people in remote imagery & Real-time speed untested & \citep{Chen2025} \\
U-ShapeNet, U-shaped extraction and fusion & RGB & Small UAV targets & Feature reuse discards information & \citep{Jiang2025} \\
Enhancement subnetwork with RetinaNet or Faster R-CNN & RGB & Objects in low light & No gain on medium-sized objects & \citep{Wu2023} \\
MFIL-FCOS with ResNet-50 & RGB & Multi-scale remote sensing targets & Struggles on highly irregular scenes & \citep{Zhang2024} \\
YOLOv5s with SPDA-C3 and Res-DHead & RGB & Dense aerial targets & Speed reported on high-end GPUs only & \citep{Sun2023} \\
YOLOv8 with hybrid MobileNetV3 stages & RGB & Vehicles & Not verified on real mobile hardware & \citep{Liu2026a} \\
ShipFPN, ResNet with spectral convolutions & RGB & Maritime vessels in satellite imagery & Cost grows with search area & \citep{Mbietieu2025} \\
Faster R-CNN with FaceNet and DBSCAN clustering & RGB & Strawberries and flowers from multiple views & Long pipeline with many stages & \citep{Zheng2023} \\
MCFAD-YOLO, YOLOv11s with DS-PAN & RGB & Small remote sensing targets & Slower than standard lightweight YOLO models & \citep{Wang2026a} \\
Qwen2.5-Omni with adaptive transformer tuning & RGB & Military targets & Little variety of real environments & \citep{Zhang2026c} \\
YOLOv5 with squeeze-and-excitation attention & RGB & \textbf{Boiler inner wall defects} & Attention modules add parameters and complicate training & \citep{Sun2022} \\
DCGAN augmentation with YOLOv5-SE & RGB & \textbf{Corroded bridge bolts} & Severity labeled by hand; no progression over time & \citep{Li2024b} \\
Bent identity CNN with fine-tuning & RGB & Aerial remote sensing targets & Long pipeline requiring complex segmentation & \citep{Lalitha2022} \\
SSD and YOLOv8 & CT & Midfacial fractures & Single hospital, limited variety & \citep{Morita2024} \\
YOLOv8-LSK with SAHI and MPDIoU loss & Radar & Oil spills & Boxes cannot map the spill outline & \citep{Zhao2026} \\
SonarDeNet with shadow-prior and kernel attention & Sonar & Underwater targets & Does not model acoustic reverberation & \citep{He2026} \\
Domain-adaptive Faster R-CNN with FPN & X-ray & Voids in spacecraft composites & Requires labels in both source and target domains & \citep{Gong2022} \\
\end{xltabular}
\endgroup

\subsection{Backbones and Transfer Learning}

Convolutional backbones remain the default feature extractor across this
literature, with transformers appearing as an auxiliary branch, an ensemble
member or a comparison baseline rather than as a replacement
\citep{Song2024,Gao2023,Zhang2026b,Bechis2026,Yamaguchi2024,Chowdhury2026,Su2024,Kharyuk2026}. Lightweight and mobile backbones are adopted
where inference cost matters \citep{Liu2026a,Mirzaei2025}, and residual or dense families
elsewhere \citep{Zhang2024,Yamaguchi2024,Su2024}.

Transfer learning is the standard remedy for limited annotation. \citet{Su2024}
fine-tunes ImageNet-pretrained VGG-16 and MobileNet and then reuses them as fixed
feature extractors, reporting faster convergence and higher accuracy than
training from scratch; \citet{Gong2022} transfers across domains to detect defects from
few samples; and \citet{Zhang2026c} adapts a pretrained multimodal model through a small
task-specific head. This matters here because a trimodal corrosion dataset of the
scale conventionally required does not exist, a point developed in
Section~\ref{sec:lr-datasets}.

Two studies in this group carry structural lessons for fusion. \citet{Chowdhury2026} compares
early, intermediate and late fusion within a single architecture and finds early
fusion strongest, then distills the result to cut computation by 60\%, although
its second stream is a tabular record rather than a sensor. \citet{Zhang2026b} weights
modalities by attention inside a YOLOv11 detector on visible and thermal pavement
imagery, and is the only study in this group to use two physically distinct
imaging sensors; it requires the pairs to be aligned in advance and loses 20
accuracy points when moved from its own data to a public benchmark.

\subsection{Interpretability and Label-Efficient Learning}

Two capabilities this research requires are established but have not been applied
to fused sensing. The first is attribution. \citet{Moradi2024} explains detector output by
hierarchical random masking without access to model internals, and \citet{Persiya2025}
attributes each input feature through SHAP values, so the contribution of an
individual channel can in principle be measured rather than asserted. The second
is learning without dense labels. \citet{Phannakham2026} grades corrosion-induced deterioration
from acoustic signals by unsupervised clustering of frequency-domain features,
requiring no annotation and no expert interpreter, though the method holds only
above roughly 4\% mass loss and was validated on five laboratory beams. \citet{Mirzaei2025}
recovers targets a few pixels across through a deliberately shallow network, and
\citet{Hafdaoui2026} reduces false alarms by pairing convolutional features with a support
vector machine. Table~\ref{tab:architectures} summarizes this group.

\begingroup
\scriptsize
\renewcommand{\arraystretch}{1.1}
\begin{xltabular}{\textwidth}{@{}X p{2.2cm} p{2.4cm} X l@{}}
\caption{Summary of the backbone, interpretability and label-efficiency studies reviewed.}\label{tab:architectures}\\
\toprule
\textbf{Method} & \textbf{Role} & \textbf{Sensing} & \textbf{Reported limitation} & \textbf{Ref.} \\
\midrule
\endfirsthead
\multicolumn{5}{@{}l}{\footnotesize\itshape Table~\thetable\ (continued)}\\
\toprule
\textbf{Method} & \textbf{Role} & \textbf{Sensing} & \textbf{Reported limitation} & \textbf{Ref.} \\
\midrule
\endhead
\midrule
\multicolumn{5}{r@{}}{\footnotesize\itshape Continued on next page}\\
\endfoot
\bottomrule
\endlastfoot
U-Net with a Transformer branch and attention gates & Segmentation backbone & Polarimetric channels, one SAR sensor & Slower per image; overfits on small fragmented targets & \citep{Song2024} \\
Parallel geometric and polarization backbones & Detection backbone & Two feature streams, one radar sensor & Polarization branch partly redundant; needs full polarimetry & \citep{Gao2023} \\
YOLOv11 with attention-based modality weighting & Detector with learned fusion & RGB and thermal & Requires pre-aligned pairs; 20-point drop across datasets & \citep{Zhang2026b} \\
Recurrent CNN with ConvLSTM over depth maps & Sequence encoder & LiDAR only & Simulation only; accuracy varies with angular velocity & \citep{Bechis2026} \\
VGG-16 with autoencoder and self-organizing map & Classifier with clustering & RGB & Classification rather than localization; 1--4\% gain & \citep{Yamaguchi2024} \\
Graph convolutional network fused with a CNN at input & Early-fusion classifier & RGB and tabular records & Second stream is not a sensor; medical domain & \citep{Chowdhury2026} \\
Pretrained VGG-16 and MobileNet fine-tuned & Transfer-learned backbone & RGB & Width regression far weaker than detection; concrete not metal & \citep{Su2024} \\
Existing backbones probed under image perturbation & Diagnostic study & None & Expensive to run and store & \citep{Kharyuk2026} \\
BODEM, hierarchical random masking & Post-hoc explanation & Detector-agnostic & Heavy to run; small test set; no human study & \citep{Moradi2024} \\
Five-layer CNN with dilation preprocessing & Small-target detector & Optical microscopy & Two classes; high-contrast laboratory colloids only & \citep{Mirzaei2025} \\
CNN features with a support vector machine & Classifier & Multispectral satellite & Reported without stated limitation & \citep{Hafdaoui2026} \\
One-dimensional CNN with SHAP attribution & Explainable classifier & Thermal and behavioral features & 197 imbalanced samples; data cannot be shared & \citep{Persiya2025} \\
Frequency-domain transforms with K-means++ clustering & Unsupervised grading & Acoustic & Valid only above roughly 4\% mass loss; five laboratory beams & \citep{Phannakham2026} \\
\end{xltabular}
\endgroup

Three findings follow. First, the architectural machinery this research needs is
mature and available, and no new detector family is required; the contribution
must lie in what the detector is given rather than in the detector itself.
Second, none of the 13 studies in Table~\ref{tab:architectures} trains on more
than two physically distinct imaging sensors, and the single study that uses two
\citep{Zhang2026b} requires the streams to be registered beforehand rather than learning the
correspondence. Third, attribution \citep{Moradi2024,Persiya2025} and label-efficient learning
\citep{Phannakham2026} have each been demonstrated on single-channel data, which makes the
modality-contribution question posed in Chapter~\ref{ch:introduction}
answerable in principle. \citet{Ding2025} has posed it across a measured
trimodal input and answered it; no reviewed study has posed it across a
contactless one.

\section{Datasets and Evaluation Practice}
\label{sec:lr-datasets}

The architectures of Section~\ref{sec:lr-architectures} are trained on data, and
the data available to corrosion research constrains what can be asked of them
more tightly than the architectures do. This section examines what the reviewed
studies train on, under what conditions that data was acquired, and how results
are reported.

\subsection{Data Availability}

No dataset identified in this review pairs visible, thermal and geometric
measurements of a corroded metal surface. Every corrosion study reviewed in
Sections~\ref{sec:lr-corrosion} to \ref{sec:lr-modalities} acquires its own
data, and none of those studies reports releasing it. The scarcity is not
peculiar to the trimodal case: \citet{Nash2022} attributes the field failures
it documents to the absence of a public corrosion dataset, so the constraint
already binds at a single modality. The search described in
Section~\ref{sec:lr-introduction} covers the publication databases rather than
the data repositories, so this is a statement about the datasets accompanying
the reviewed literature and not a proof that no such collection exists anywhere;
a repository-level search forms part of the acquisition stage in
Chapter~\ref{ch:methodology}.
The three studies that do sense all three modalities concurrently
\citep{Tran2025,Yang2024,Tran2024} release driving scenes and calibration sequences, which carry no
corrosion and no annotation of material loss. Of the corrosion studies that fuse modalities,
only \citet{Ding2025} exceeds two channels, and it builds its own dataset for a
contact-based sensor set. The contactless pairs stop at two, as
Section~\ref{sec:lr-fusion} sets out, and neither releases its data.

This is the practical form of the gap set out in Section~\ref{sec:lr-gap}. A
trimodal corrosion model cannot be trained by assembling existing public data,
because the modality combination and the target do not co-occur in any released
collection. Acquisition is therefore part of the problem rather than a
precondition for it, and the methodology in Chapter~\ref{ch:methodology} has to
answer for it.

\subsection{Acquisition Conditions}

Where data is collected, it is usually collected under conditions more
favorable than the intended deployment. Laboratory specimens stand in for field
surfaces in several of the most directly relevant studies: \citet{Maru2025} evaluates
section loss on fabricated steel components and states that real corrosion
texture is absent, \citet{Asad2026} trains on synthetic and laboratory shapes, and
\citet{Phannakham2026} validates on five laboratory beams. Where field data is used, the
acquisition window is narrow. \citet{Guedes2025} requires the dry season to fly, \citet{Yamashita2025}
reports its method unusable for several days after rain because residual
moisture masks the thermal signature, and \citet{Lim2021} requires halogen lamps to
excite the surface, an arrangement that does not extend to large structures.
The conditions inside a
collection matter as much as the setting it was gathered in: \citet{Yao2019}
report that their classifier generalizes poorly to images taken under weak
illumination, blur and shade, and attribute this to those conditions being
under-represented in the training set rather than to the architecture.
\citet{Jiang2023} is the one reviewed study to quantify the effect rather than
note it. Working inside dim steel box girders, it runs photography experiments
across shooting distance, focal length and view angle, and derives an equation
relating detected pixel area to physical area because the two are not related by
a fixed factor. The result bears on more than data collection: it means that
even affected area, the quantity appearance-based severity grading reports, is a
function of where the camera stood, and it is recovered from a database of 300
images after augmentation.
Section~\ref{sec:lr-modalities} established that thermal contrast depends on an
excitation source; this is where that dependence becomes a data-collection
constraint rather than a sensing one.

One pattern across these constraints bears on what fusion can be expected to
deliver. The conditions that degrade one channel are not always independent of
those that degrade another. Surface moisture is the clearest case: it masks the
thermal signature for several days after rain \citep{Yamashita2025}, and wet
surfaces are named among the conditions that obscure damage features from an
appearance model \citep{Maharjan2025}. A fusion argument implicitly assumes that
where one modality is compromised another is not, and shared failure modes
weaken that assumption rather than defeat it. The consequence is for acquisition
rather than for architecture: Chapter~\ref{ch:methodology} has to state the
conditions under which data is and is not collected, instead of treating
coverage of the sensor set as sufficient on its own.

\subsection{Annotation and Sample Size}

Annotation is the second constraint. \citet{Li2024b} assigns its three corrosion
severity classes by hand and notes that no progression over time is captured;
\citet{Gong2022} requires labeled examples in both the source and target domains for its
transfer to work. Sample sizes are correspondingly small: \citet{Persiya2025} reports 197
imbalanced samples and states that its thermal data cannot be shared, and
\citet{Morita2024} draws from a single institution with limited variety.

Two responses appear. The first is synthesis: \citet{Li2024b} generates additional
corroded bolt images with a generative adversarial network, and \citet{Sun2022} applies
conventional augmentation to enlarge a defect set. The second is to remove the
labeling requirement altogether, as \citet{Phannakham2026} does by clustering frequency-domain
features without annotation, at the cost of a detection floor near 4\% mass
loss. Both remain untested on multimodal input.
Table~\ref{tab:data} groups these constraints.

\begin{table}[htbp]
\centering
\caption{Data and evaluation constraints reported across the reviewed literature.}
\label{tab:data}
\scriptsize
\renewcommand{\arraystretch}{1.15}
\begin{tabularx}{\textwidth}{@{}p{3.2cm} X p{3.4cm}@{}}
\toprule
\textbf{Constraint} & \textbf{How it appears in the literature} & \textbf{Ref.} \\
\midrule
No public trimodal corrosion data & Every corrosion study reviewed acquires its own data; the three studies that sense all three modalities release driving or calibration data & \citep{Tran2025,Yang2024,Tran2024} \\
Laboratory acquisition & Specimens, benches or simulations stand in for field surfaces, so real corrosion texture and clutter are absent & \citep{Maru2025,Asad2026,Phannakham2026,Mirzaei2025,Zhang2026c,Bechis2026} \\
Environmental windows & Acquisition is restricted to conditions the sensor tolerates: dry season for flight, several rain-free days, or an external heat source & \citep{Guedes2025,Yamashita2025,Lim2021} \\
Annotation cost & Severity classes assigned by hand, or labels required in both source and target domains & \citep{Li2024b,Gong2022} \\
Small and imbalanced samples & Tens to low hundreds of examples, unevenly distributed across classes, sometimes unshareable & \citep{Persiya2025,Phannakham2026,Morita2024} \\
Synthesis as a remedy & Generative augmentation and conventional augmentation stand in for scarce real examples & \citep{Li2024b,Sun2022} \\
Untested generalization & Performance reported on one collection; where a second is tried, accuracy falls sharply & \citep{Zhang2026b,Morita2024,Zhang2026c} \\
\bottomrule
\end{tabularx}
\end{table}

\subsection{Evaluation Practice}

Detection is reported consistently. Mean average precision, intersection over
union, precision, recall and the F1 score recur across the detection and fusion
literature, so models can be compared against one another on a shared basis.
Severity is not, for the reasons set out in Section~\ref{sec:lr-corrosion}:
almost half of the corrosion studies reviewed do not quantify material loss at
all, and those that do adopt measures that cannot be converted into one
another. There is
consequently a benchmark for finding corrosion and none for measuring it.

Two consequences follow. First, no shared severity benchmark exists against
which a depth-recovering method could be scored, so a study proposing one must
also state the measure it is scored by and supply its own reference
measurements. This research adopts maximum pit depth within a detected region,
reported as mean absolute error against an independently instrumented
reference, and Chapter~\ref{ch:methodology} specifies how that reference is
obtained. In the absence of a shared benchmark the comparator is internal: the
depth error of each ablation configuration is read against the depth error of
the appearance-only configuration, so the geometric channel is scored on what it
adds rather than against an external threshold that does not exist. Second, cross-dataset generalization is
rarely tested, and the one reviewed study that tests it directly loses 20
accuracy points when moved from its own collection to a public benchmark
\citep{Zhang2026b}. Reported performance in this literature is therefore a statement about
a particular collection rather than about a method, which is the reason the
ablation proposed in Chapter~\ref{ch:introduction} is designed to hold data,
splits and training conditions fixed while varying only the sensing
configuration.

\section{Research Gap}
\label{sec:lr-gap}

The literature reviewed in this chapter supports a chain of reasoning that
terminates in a specific absence. Corrosion is an electrochemical process whose
most hazardous form is localized, and its hazard depends on depth as well as extent
\citep{Qin2026,Liu2026,Wang2021,Akhlaghi2023}. Inspection practice measures depth
with contact instruments that require access to the surface
\citep{Huang2026a,Yang2025}; image-based severity estimation is mostly confined
to the area it can see \citep{Alvani2026}, and the studies that recover depth
optically do so on prepared specimens under controlled conditions
\citep{Egodawela2025,Doshvarpassand2022,Haitz2022,Maru2025}. \citet{Alvani2026} puts the
division in its sharpest form: frameworks combining accurate detection,
localization and explicit quantitative severity estimation remain scarce, and
severity across most studies is either omitted, inferred heuristically from
global features, or dependent on specialized hardware. Those three cases
describe the reviewed literature closely. The fourth, a quantitative severity
produced jointly with the detection and without a contact instrument, is the one
this research targets, and the paper that names the division occupies the second
case itself. No single modality closes this. Visible imaging
records the surface and carries no depth information \citep{Cheng2025}; infrared
thermography reveals subsurface loss but depends on thermal contrast and, in the
one steel-corrosion demonstration reviewed, on an external heat source
\citep{Lim2021,Doshvarpassand2022}; geometric sensing measures shape without distinguishing corrosion
product from intact metal \citep{Maru2025,PimenteldeFigueiredo2025}.

Fusion is the established answer to a limitation of this shape, and
Section~\ref{sec:lr-fusion} finds it mature across the application domains
represented in this review and largely absent from defect detection. The tally
there is stark: 3 of the 38 tabulated studies sense visible, thermal and
geometric data concurrently, and each of the 3 directs the combination at
driving scenes or at calibrating the rig itself \citep{Tran2025,Yang2024,Tran2024}. The
corrosion end of the same table is bimodal throughout
\citep{Lim2021,PimenteldeFigueiredo2025}. Within the journal
literature reviewed here the position is simpler than a tally of near-misses: no
study fuses visible, thermal and geometric sensing for a corrosion decision at
all, whether the combination is learned or assembled by hand, and none recovers
a depth measurement from a fused representation.

One conference paper does reach part of this. \citet{Ding2025} learns across
three modalities for a corrosion decision and reports what each contributes. It
lies outside the corpus this review tabulates and the gap is not argued from it;
it is set out because a stated absence that ignores the nearest attempt is a
weaker claim than one that names it. Two things separate it from what is
proposed here. Its dominant modality is ultrasonic and requires transducers
mounted on the surface, so the sensor set is not contactless and does not
address the access constraint; and its output is an ordinal severity class
rather than a measured depth. It therefore reaches neither the sensing nor the
output this research targets. The absence also registers from inside the literature
and not only in a tally taken across it: \citet{Imran2025} close a
single-modality study by naming the fusion of optical, thermal and acoustic
sensing as an open direction, which places the combination among the things not
yet attempted rather than among those tried and abandoned.

The absence is best stated precisely, as a property of the sensor set and the
output rather than as a claim about the completeness of the search. No study
reviewed here learns a joint representation across visible, thermal and
geometric channels for a corrosion decision, and the nearest attempt outside
that corpus reaches neither the contactless sensing nor the measured output, for
the reasons given above. Elsewhere, where a third
channel appears in an inspection setting it is inferred from imagery rather than
measured by a dedicated sensor \citep{Cao2026,Nooralishahi2022}, and where three
modalities are genuinely acquired the streams are combined at the reporting
stage through layer-by-layer comparison \citep{Guedes2025} rather than learned
from data. The proposition this research tests is whether
thermography can occupy, contactlessly, the subsurface-sensing role that
inspection practice currently fills with a probe placed on the surface
\citep{Huang2026a,Yang2025}. The review does not settle this.
Every steel-corrosion thermography result surveyed in
Section~\ref{sec:lr-modalities} is excited rather than passive
\citep{Lim2021,Doshvarpassand2022,Elizalde2023}, contrast attenuates with
defect depth, and the excitation arrangement that \citet{Lim2021} uses does not
extend to inspection standoff. The excitation regime is therefore a
methodological question this research must answer rather than a property it can
assume, and it is treated as such in Chapter~\ref{ch:methodology}. Two distinctions matter. The first is access.
Section~\ref{sec:lr-corrosion-detection} established that the regions most
exposed to attack are frequently the least reachable, so a method requiring a
probe to be placed on the surface inherits the constraint that optical
inspection exists to remove; a contactless set can be carried past a structure
rather than applied to it. The second is that a fixed rule cannot represent the
interaction this problem turns on, namely that a region is
corroded when a discoloration, a thermal anomaly and a depression coincide, and
that the evidence from each channel is conditional on the other two.

Three obstacles explain the absence, and each is documented above. The first is
spatial correspondence. Section~\ref{sec:lr-fusion} found 9 of the tabulated
fusion studies naming alignment, calibration or registration failure among their
limitations, and the difficulty is compounded for thermal sensors because a
printed calibration target has no thermal contrast \citep{Yang2024}. Two studies
do calibrate all three modalities without a target \citep{Yang2024,Tran2024}, so
the correspondence is recoverable; neither carries the aligned streams forward
into a detection model, and that is where the gap lies rather than in the
calibration itself. The second is
data: no released collection pairs the three modalities with a corroded metal
surface, so acquisition is part of the problem rather than a precondition for it
(Section~\ref{sec:lr-datasets}). The third is evidence: fusion is usually compared against a
single unimodal baseline rather than across the sensor set, so the contribution
of an individual modality is asserted rather than measured. Four studies in the
reviewed corpus are exceptions
\citep{Hou2024,Haitz2022,Maru2025,Rashidi2026}, and none of them isolates a
ranging channel on a corrosion task. The only corrosion study to vary a sensor
set across three modalities is \citet{Ding2025}, outside that corpus, where the
geometric channel accounts for 2.9\% of classification accuracy against 9.2\%
for the contact ultrasonic channel. The figure does not transfer: the decision
there is an ordinal class rather than a depth measurement, and the geometric
channel is not the only source of the quantity being predicted. It is reason to
treat the contribution of a ranging sensor as something to be measured rather
than assumed, which is what Objective~3 does.

What the literature does not lack is machinery. Section~\ref{sec:lr-architectures}
establishes that detectors, convolutional backbones, transfer learning,
attribution and label-efficient training are all mature and available
\citep{Srivastava2021,Su2024,Moradi2024,Phannakham2026}. The contribution this research proposes therefore lies not in
a new architecture but in what the architecture is given and in what is then
measured: a registration pipeline that brings visible, thermal and geometric
measurements of a metal surface into a common frame; a fusion-based detector
that learns a joint representation from the aligned result and reports, for each
detection, the maximum pit depth within it; and an ablation that isolates each
modality's contribution to both outputs under fixed data, splits and training
conditions.

The depth output changes what the ablation is asking.
Section~\ref{sec:lr-corrosion} established that depth matters to the hazard alongside extent while the
surface understates it, and Section~\ref{sec:lr-corrosion-detection} that
image-based severity estimation is routinely confined to affected area
\citep{Alvani2026}. That confinement is a convention rather than a
physical limit, but the evidence that appearance can replace range is weaker
than it first appears. \citet{Egodawela2025} regresses depth to an RMSE of
0.2829~mm from multispectral reflectance, and is the figure most often cited for
appearance-based depth; the defects are machined flat-bottom holes drilled into
the face the camera observes, so the quantity recovered is the geometry of an
open cavity rather than loss beneath an intact surface, and acquisition is gated
by a thermocouple in contact with the specimen.
\citet{Doshvarpassand2022} matches thermal contrast curves against an analytic
model for defects machined into the hidden face, which is genuinely subsurface,
but reports no error figure and states that accuracy falls as the defect
aspect ratio falls. Neither is a like-for-like baseline for a camera viewing a
corroded surface. The appearance-only arm of the ablation is therefore a
necessary control rather than a strong competitor, and the literature does not
establish how well it will perform. The
question Objective~3 puts to the geometric channel is therefore whether direct
range measurement improves on what appearance alone already achieves, by a
margin that justifies carrying the sensor. A negative answer is a reportable
result. These are the three objectives stated in
Chapter~\ref{ch:introduction}, and the methodology in
Chapter~\ref{ch:methodology} sets out how each is to be met.
